\documentclass[../cuaderno6.tex]{subfiles}

\begin{document}

\subsection{Disco horizontal}

En el caso del disco medido en posición horizontal, se encuentra un resultado bastante interesante al ver la magnetización remanente y el campo coercitivo, visibles en la figuras \ref{fig:H_c_segun_H_a_Disco_horizontal_Fa0_17_09} y \ref{fig:M_r_segun_H_a_Disco_horizontal_Fa0_17_09}, que es que para campos bajos, predomina la estructura microscópica del acero. Se encuentra una variación según el ángulo no esperada para esta geometría, ya que el grano del material hace que este responda de forma diferente según su orientación. La magnetización máxima, visible en la figura \ref{fig:M_max_segun_H_a_Disco_horizontal_Fa0_17_09} no presenta variación alguna según el ángulo.

\begin{figure}[H]
    \centering
    \includegraphics[width=0.5\textwidth]{../img/Fa0/Disco horizontal/H_c_segun_H_a_Disco horizontal_Fa0_17_09.png}
    \caption{El campo coercitivo $H_c$ en función del logaritmo del campo aplicado $H_{max}$ para cada ángulo.}
    \label{fig:H_c_segun_H_a_Disco_horizontal_Fa0_17_09}
\end{figure}
\begin{figure}[h]
    \centering
    \includegraphics[width=0.5\textwidth]{../img/Fa0/Disco horizontal/M_r_segun_H_a_Disco horizontal_Fa0_17_09.png}
    \caption{Magnetización remanente $M_r$ en función del campo aplicado $H_{max}$.}
    \label{fig:M_r_segun_H_a_Disco_horizontal_Fa0_17_09}
\end{figure}
\begin{figure}[h]
    \centering
    \includegraphics[width=0.5\textwidth]{../img/Fa0/Disco horizontal/M_max_segun_H_a_Disco horizontal_Fa0_17_09.png}
    \caption{Magnetización máxima $M_{max}$ en función del campo aplicado $H_{max}$.}
    \label{fig:M_max_segun_H_a_Disco_horizontal_Fa0_17_09}
\end{figure}

La variación del resultado del ajuste para el campo coercitivo resulto mínima entre las obtenidas anteriormente. La variación en el factor de desmagnetización es básicamente nula, solo aparece a partir de la cuarta cifra. Es tiene sentido, girar un disco sobre su eje de rotación no debería inducir ningún cambio en su geometría, ya que es una de sus simetrías. Vemos que esta se respeta.

\begin{figure}[H]
    \centering
     \includegraphics[width=0.5\textwidth]{../img/Fa0/Disco horizontal/H0_c_angulo_minors_Disco horizontal.png}
    \caption{Gráfico de la pendiente del ajuste $H_c^0$ en función del ángulo.}
    \label{fig:H0_c_angulo_minors_Disco_horizontal}
\end{figure}
\begin{figure}[H]
    \centering
    \includegraphics[width=0.5\textwidth]{../img/Fa0/Disco horizontal/N_d_angulo_major_Disco horizontal.png}
    \caption{Factor de desmagnetización $N_d$ en función del ángulo.}
    \label{fig:N_d_angulo_major_Disco}
\end{figure}

\end{document}